\section{Mode 2: Program Validation and Execution Mode}

The second operating mode of the digital twin is the program validation and execution mode. Unlike monitoring mode, where the physical robot drives the simulated robot, program mode starts from a saved robot program. The program is first tested inside the digital twin environment, and only after successful validation can it be sent to the physical controllers.

This mode connects the teaching workflow to the real robot cell. The teach panel creates a YAML program that represents the operator's intended sequence of arm waypoints and gripper commands. Program mode then checks that sequence through MoveIt and the Gazebo twin, stores the validation result, and enforces a validation-before-execution rule before any real motion is requested.

\begin{figure}[H]
    \centering
    \resizebox{\textwidth}{!}{%
    \begin{tikzpicture}[
        font=\small,
        node distance=0.85cm and 1.1cm,
        data/.style={draw, rounded corners=2pt, align=center, minimum width=2.75cm, minimum height=0.85cm, fill=gray!8},
        logic/.style={draw, rounded corners=2pt, align=center, minimum width=3.2cm, minimum height=0.9cm, fill=orange!10},
        twin/.style={draw, rounded corners=2pt, align=center, minimum width=2.75cm, minimum height=0.85cm, fill=green!8},
        real/.style={draw, rounded corners=2pt, align=center, minimum width=2.8cm, minimum height=0.85cm, fill=blue!8},
        topic/.style={draw, rounded corners=2pt, align=center, minimum width=2.75cm, minimum height=0.8cm, fill=purple!8},
        arrow/.style={-Latex, thick},
        feedback/.style={-Latex, thick, dashed}
    ]
        \node[data] (yaml) {Saved YAML\\program};
        \node[logic, right=of yaml] (simrunner) {\texttt{program\_sim\_runner.py}\\validation runner};
        \node[twin, right=of simrunner] (moveit) {MoveIt action\\\texttt{/move\_action}};
        \node[twin, right=of moveit] (gazebo) {Gazebo twin\\execution};
        \node[topic, above=of simrunner] (validate) {\texttt{/dt/program\_validate}};

        \node[data, below=1.1cm of simrunner] (validated) {Validated YAML\\saved trajectories};
        \node[logic, right=of validated] (realrunner) {\texttt{program\_real\_executor.py}\\execution gate};
        \node[real, right=of realrunner] (controllers) {Real robot\\trajectory controllers};

        \node[topic, below=0.9cm of validated] (execute) {\texttt{/dt/program\_execute}};
        \node[topic, below=0.9cm of realrunner] (status) {\texttt{/dt/program\_status}\\\texttt{/dt/program\_progress}};
        \node[topic, below=0.9cm of controllers] (estop) {\texttt{/dt/estop}\\stop request};

        \draw[arrow] (validate) -- (simrunner);
        \draw[arrow] (yaml) -- (simrunner);
        \draw[arrow] (simrunner) -- (moveit);
        \draw[arrow] (moveit) -- (gazebo);
        \draw[feedback] (simrunner) -- (validated);

        \draw[arrow] (execute) -- (realrunner);
        \draw[arrow] (validated) -- (realrunner);
        \draw[arrow] (realrunner) -- (controllers);
        \draw[arrow] (estop) -- (realrunner);

        \draw[feedback] (realrunner) -- (status);

        \begin{scope}[on background layer]
            \node[draw, dashed, rounded corners=4pt, fit=(yaml)(validate)(simrunner)(moveit)(gazebo), inner sep=0.28cm, label=above:Simulation validation path] {};
            \node[draw, dashed, rounded corners=4pt, fit=(execute)(validated)(realrunner)(controllers)(status)(estop), inner sep=0.28cm, label=below:Real execution path] {};
        \end{scope}
    \end{tikzpicture}%
    }
    \caption{Program validation and execution mode workflow.}
    \label{fig:dt_program_mode}
\end{figure}

\subsection{Purpose of Program Mode}

The purpose of program mode is to reduce the risk of moving from an operator-taught program to real robot execution. A saved waypoint is only a target recorded by the user; it does not automatically prove that the robot can reach the pose, that the joint vector is valid for the selected robot, or that the sequence can be executed by the available controllers. Program mode adds an intermediate stage where the digital twin checks the program before the real cell is allowed to run it.

This is especially important in the dual-arm setup because the program may contain steps for both the ABB IRB120 and the AR4 MK3, as well as gripper commands. The digital twin provides a controlled environment where these steps can be processed in order, visualized, and rejected if the planning or execution pipeline fails.

\subsection{Validation Request and Program Loading}

Validation begins when the teach panel publishes the program name on \texttt{/dt/program\_validate}. The node \texttt{program\_sim\_runner.py} receives this request and loads the corresponding YAML file from \texttt{\textasciitilde/.ros/dt\_programs}. The runner checks that the file exists, that it is a valid YAML map, and that it contains a non-empty ordered list of waypoints.

Only one validation run is allowed at a time. If a validation request arrives while another program is already being processed, the new request is rejected and a status message is published. This prevents two programs from trying to control the same twin planning and execution pipeline at the same time.

\subsection{Simulation Validation Using MoveIt}

After loading the program, the validation runner processes each step in sequence. For an arm waypoint, it checks the selected robot, verifies that the stored joint vector has the correct length, and builds a MoveIt \texttt{MoveGroup} goal for the correct planning group. The ABB arm uses the \texttt{irb120\_arm} group, while the AR4 arm uses the \texttt{ar\_manipulator} group.

Gripper steps are validated through the same general mechanism. The requested \texttt{open} or \texttt{closed} state is converted into the appropriate gripper joint target for the selected robot, then sent to MoveIt as a planning and execution request for the gripper group. This keeps motion steps and gripper steps inside one ordered validation sequence.

The MoveIt request is sent through the \texttt{/move\_action} action interface. In this mode, the execution target is the simulation environment, not the physical robot. A successful result means that MoveIt accepted the goal, produced a valid planned trajectory, and executed it in the Gazebo twin. During this process, the runner publishes progress and textual status feedback on \texttt{/dt/program\_progress} and \texttt{/dt/program\_status}.

\subsection{Saving Validation Results}

When a waypoint succeeds, the runner marks that step as validated and stores the planned trajectory inside the same YAML program. The saved trajectory contains the joint names and trajectory points returned by MoveIt. If a step fails, the runner stops the validation sequence, marks the failed step as unvalidated, and leaves the program-level validation flag false.

At the end of a successful validation run, the program is saved again with \texttt{validated: true}. This updated YAML file becomes the formal output of the twin validation stage. The file no longer represents only the operator's taught intention; it also records that the digital twin planning pipeline was able to execute the sequence in simulation.

\subsection{Validation Gate Before Real Execution}

The validation flag is used as a gate before real execution. The teach panel checks the saved program before publishing an execution request, and the real execution runner performs the same check again after loading the file. If \texttt{validated} is not true, \texttt{program\_real\_executor.py} rejects the request and publishes an error status instead of commanding the controllers.

This duplication is deliberate. The interface prevents common operator mistakes, while the backend runner enforces the rule even if an execution request is published manually from another ROS node. Therefore, validation is not treated as a visual indication only; it becomes part of the execution contract between the digital twin and the physical system.

\subsection{Real-Hardware Program Execution}

Real execution begins when a validated program name is published on \texttt{/dt/program\_execute}. The node \texttt{program\_real\_executor.py} loads the YAML file, checks the program-level validation flag, and then executes the ordered list of steps on the physical controllers.

For arm waypoints, the executor commands the corresponding \texttt{FollowJointTrajectory} action server. The ABB arm is sent to \texttt{/irb120\_controller/follow\_joint\_trajectory}, while the AR4 arm is sent to \texttt{/ar4\_trajectory\_controller/follow\_joint\_trajectory}. The executor reads the latest real joint states and builds a trajectory from the current measured position to the saved target joint vector. When configured to use saved trajectories, it can instead replay the planned trajectory stored during validation, after checking that the saved joint names match the target controller.

For gripper commands, the same runner sends trajectory goals to the ABB or AR4 gripper controller. The gripper step is therefore executed in its proper place in the program sequence, rather than being handled as a separate manual command outside the saved program.

\subsection{Stop Handling and Status Feedback}

Program mode also provides runtime feedback during validation and execution. Both runners publish human-readable status messages on \texttt{/dt/program\_status} and normalized progress values on \texttt{/dt/program\_progress}. These topics allow the teach panel or another operator interface to show whether the system is idle, validating, executing, complete, or blocked by an error.

The real execution runner subscribes to \texttt{/dt/estop}. A message containing \texttt{stop}, \texttt{estop}, or \texttt{halt} sets a stop flag and attempts to cancel the active trajectory goal. In this thesis, this mechanism should be understood as a program-level stop and cancellation feature. It improves operator control over the program runner, but it is not a certified hardware safety layer. Physical emergency stops and controller-level safety functions remain separate from the digital twin software.

\subsection{Contribution to the Digital Twin Task}

Program mode completes the digital twin workflow by turning the twin into a pre-execution validation environment. Monitoring mode keeps the virtual cell synchronized with real motion, while program mode uses the same robot models, controllers, and planning environment to test saved programs before they reach the real hardware.

This directly supports the collaborative assembly task. The operator can teach a sequence, validate it in the twin, inspect the feedback, and then execute only the approved program on the physical cell. The result is a clearer separation between teaching, validation, and execution, which makes the system easier to operate and reduces the chance that an untested program is sent directly to the robots.
